%% draft_fatou_hopf.tex
%% Draft research article — November 2026
%% Authors: V. Álvarez-Aparicio, J.M. García-Calcines,
%%          L.J. Hernández-Paricio, M.T. Rivas-Rodríguez
%%
%% Compile: pdflatex draft_fatou_hopf.tex (twice)

\documentclass[12pt]{article}

\usepackage{amsmath, amssymb, amsthm}
\usepackage{mathtools}
\usepackage{geometry}
\geometry{a4paper, margin=2.5cm}
\usepackage{hyperref}
\usepackage{enumitem}
\usepackage{microtype}

%% ---- Theorem environments ------------------------------------------------
\newtheorem{theorem}{Theorem}[section]
\newtheorem{proposition}[theorem]{Proposition}
\newtheorem{lemma}[theorem]{Lemma}
\newtheorem{corollary}[theorem]{Corollary}
\newtheorem{conjecture}[theorem]{Conjecture}
\theoremstyle{definition}
\newtheorem{definition}[theorem]{Definition}
\newtheorem{example}[theorem]{Example}
\newtheorem{remark}[theorem]{Remark}
\newtheorem{algorithm_desc}[theorem]{Algorithm}

%% ---- Macros ---------------------------------------------------------------
\newcommand{\C}{\mathbb{C}}
\newcommand{\R}{\mathbb{R}}
\newcommand{\Z}{\mathbb{Z}}
\newcommand{\N}{\mathbb{N}}
\newcommand{\Q}{\mathbb{Q}}
\newcommand{\Chat}{\hat{\C}}
\newcommand{\Pone}{\mathbf{P}^1(\C)}
\newcommand{\Poneplus}{\mathbf{P}^{1+}(\C)}
\newcommand{\Sthree}{S^3}
\newcommand{\fsh}{f^{\#}}            % spherical derivative
\newcommand{\Lfs}{L_{f^S}(\phi^S)}
\newcommand{\Lfsrk}{L_{f^S}^{[r,k]}(\phi^S)}
\newcommand{\RS}{R^S}
\newcommand{\RP}{R^P}
\newcommand{\Bf}{B_f}
\newcommand{\Jf}{J(f)}
\newcommand{\Ff}{F(f)}
\DeclareMathOperator{\Arg}{Arg}
\DeclareMathOperator{\rot}{rot}

%% ---------------------------------------------------------------------------
\title{\textbf{Classifying Fatou Components of Rational Maps\\
        via the Hopf Fibration and Lyapunov Functions}\\[6pt]
       {\large (Working draft --- November 2026)}}

\author{V.~\'Alvarez-Aparicio\thanks{Departamento de Matem\'aticas y
Computaci\'on, Universidad de La Rioja, Logro\~no, Spain.} \and
J.\,M.~Garc\'ia-Calcines\thanks{Departamento de Matem\'aticas,
Estad\'istica e I.O., Universidad de La Laguna, Spain.} \and
L.\,J.~Hern\'andez-Paricio\footnotemark[1] \and
M.\,T.~Rivas-Rodr\'iguez\footnotemark[1]}

\date{November 2026}

\begin{document}

\maketitle
\begin{abstract}
We extend the Lyapunov-based framework for rational self-maps of the Hopf
fibration $S^3 \to \Pone$ developed in \cite{Alvarez2023} to a full
classification of Fatou components.  In that earlier work, basins of
\emph{attracting} cycles were characterised by the condition that the
Lyapunov constant $c = L_{f^S}(\phi^S) < 1$, while the Julia set was
identified as the complement; no other Fatou components appeared in the
considered examples.  By Sullivan's No-Wandering-Domain theorem the full
classification of periodic Fatou components of a rational map comprises
five types: attracting/superattracting basins, parabolic basins, Siegel
discs, and Herman rings.

The main contributions of the present paper are:
\begin{enumerate}[label=(\roman*)]
  \item A \emph{convergence-rate criterion} that distinguishes parabolic
        basins from attracting basins and from the Julia set: for orbits
        tending to a parabolic cycle the Lyapunov approximation sequence
        $L_{f^S}^{[0,N]}(\phi^S)$ converges to~$1$ with polynomial rate
        $O(N^{-1/m})$, where $m$ is the parabolic multiplicity, whereas
        for attracting basins the convergence is geometric.
  \item A \emph{rotation-number algorithm on $S^3$} that detects and
        measures the irrational rotation in Siegel discs and Herman rings,
        working entirely within the normalised homogeneous coordinates of
        the Hopf endomorphism so as to avoid overflow, underflow, and the
        indeterminacy at the super-zero point $\tilde{0}$.
  \item A unified \emph{classification algorithm} that, given a rational
        map $f$ represented as an r-irreducible homogeneous pair $(F,G)$,
        assigns to each orbit on a grid of $S^3$ one of the labels:
        \textsc{attracting}, \textsc{superattracting}, \textsc{parabolic},
        \textsc{siegel}, \textsc{herman}, or \textsc{julia}.
\end{enumerate}
We illustrate the algorithms on several families of rational maps,
including Newton's method applied to cubic polynomials with parameter-
dependent cyclic behaviour, and quadratic maps with Siegel fixed points.
The algorithms are implemented in the Julia programming language as an
extension of the \texttt{RiemannBasins.jl} package \cite{RiemannBasins}.
\end{abstract}

\tableofcontents

%% ===========================================================================
\section{Introduction}
\label{sec:intro}
%% ===========================================================================

The iteration of rational maps $f: \Chat \to \Chat$ on the Riemann sphere
is one of the richest areas of one-dimensional complex dynamics.  By the
No-Wandering-Domain theorem of Sullivan \cite{Sullivan1983}, every periodic
Fatou component of a rational map of degree $d \ge 2$ belongs to exactly
one of five types: an attracting basin, a superattracting basin, a
parabolic basin, a Siegel disc, or a Herman ring \cite{Milnor2006,
Beardon1991}.  While the theoretical classification has been understood
since the 1980s, its \emph{computational} detection remains subtle:
standard orbit algorithms converge to attracting cycles with no difficulty
but give ambiguous output near parabolic cycles (slow convergence) or
inside rotation domains (quasi-periodic, non-convergent orbits).

A key technical difficulty is numerical: iterating a rational map
$f = A/B$ in affine coordinates leads to overflow when orbits pass near a
pole of $f$, to underflow when orbits approach the attracting fixed point
at $\infty$, and to outright indeterminacy when both $A(z_n) \to 0$ and
$B(z_n) \to 0$ simultaneously --- the super-zero point $[0:0]$.  In
\cite{Alvarez2023} the present authors resolved all three issues by
representing $f$ as a rational self-map $\RS: S^3 \to S^3$ of the Hopf
fibration $q: S^3 \to \Pone \cong S^2$, using normalised homogeneous
coordinates on the unit sphere $S^3 \subset \C^2$.  The spherical
derivative $\fsh: \Chat \to \R_+$ was computed entirely from bihomogeneous
polynomials $(U,V)$ continuous on all of $\C^2$ (including $(0,0)$), and
Lyapunov constants $c_i = L_{f^S}(\phi^S)|_{B_f(C_i)}$ were shown to
classify attracting cycles and their basins.

The present paper addresses the missing Fatou components.  The key
observation is as follows.  For a point $y$ in the \emph{attracting} basin
of a $q_i$-cycle $C_i$ with Lyapunov constant $c_i < 1$, the difference
between consecutive Lyapunov approximations satisfies (see
\cite[eq.~(29)]{Alvarez2023})
\[
  |L_{f^S}^{[0,k-1]}(\phi^S)(y) - L_{f^S}^{[0,k]}(\phi^S)(y)| \leq
  10^{-C}
\]
for $k$ large enough, with the gap decreasing \emph{geometrically} at rate
$c_i < 1$.  For a point in a \emph{parabolic} basin, by contrast, the
analogous sequence converges to $1$ but the gap decreases only
polynomially, at rate $k^{-1/m-1}$ where $m$ is the parabolic multiplicity
(the number of attracting petals of the fixed point of $f^{q_i}$).  For a
point in a \emph{Siegel disc or Herman ring} the orbit is quasi-periodic
and the sequence of Lyapunov approximations oscillates around $1$ without
converging.

These three regimes --- geometric decay below $1$, polynomial approach
to $1$, and oscillation around $1$ --- form the basis of our classification
algorithm.

The paper is organised as follows.  Section~\ref{sec:prelim} recalls the
Hopf fibration framework and fixes notation; we refer to \cite{Alvarez2023}
for detailed proofs and defer to \cite{Milnor2006} and \cite{Beardon1991}
for standard complex dynamics background.  Section~\ref{sec:sullivan}
reviews Sullivan's classification and recalls the relevant properties of
each Fatou type.  Section~\ref{sec:lyapunov} develops the Lyapunov-based
criteria for attracting and parabolic basins.  Section~\ref{sec:rotation}
introduces the rotation-number algorithm for Siegel discs and Herman rings
on $S^3$.  Section~\ref{sec:algorithm} presents the unified classification
algorithm.  Section~\ref{sec:examples} illustrates the method on several
families.  Section~\ref{sec:conclusions} discusses conclusions and future
directions.

%% ===========================================================================
\section{Preliminaries: the Hopf fibration framework}
\label{sec:prelim}
%% ===========================================================================

We recall briefly the setup of \cite{Alvarez2023}, to which we refer for
proofs.

\subsection{Homogeneous pairs and the Hopf endomorphism}

Let $f = A/B$ be a rational map of degree $d = \max\{\deg A, \deg B\}$.
The \emph{homogeneous pair} $(F,G) \in \C_h[z,t]\times\C_h[z,t]$
associated to $(A,B)$ is the pair of degree-$d$ homogeneous polynomials
\[
  F(z,t) = t^d A(z/t), \qquad G(z,t) = t^d B(z/t).
\]
When $(F,G)$ is r-irreducible (i.e.~irreducible as a pair with $G \ne 0$,
see \cite[Def.~12]{Alvarez2023}), there are well-defined induced maps:
\[
  R:\C^2\setminus\{(0,0)\}\to\C^2\setminus\{(0,0)\},\quad
  R(z,t) = (F(z,t),G(z,t)),
\]
\[
  \RS:\Sthree\to\Sthree,\quad
  \RS(z,t) = \frac{(F(z,t),G(z,t))}{\|(F(z,t),G(z,t))\|},
\]
\[
  \RP:\Pone\to\Pone,\quad \RP[z:t]=[F(z,t):G(z,t)],
\]
fitting into the commutative diagram
$\C^2\setminus\{0\}\xrightarrow{R}\C^2\setminus\{0\}$,
$\Sthree\xrightarrow{\RS}\Sthree$, $\Pone\xrightarrow{\RP}\Pone$,
where $p:\C^2\setminus\{0\}\to\Sthree$ is normalisation and
$q:\Sthree\to\Pone$ is the Hopf fibration $q(z,t) = [z:t]$.

\subsection{Normalisation and overflow prevention}

The \emph{normalised homogeneous pair} $[F,G]$ is defined by
\[
  [F,G] := \frac{(F,G)}{\|(F,G)\|},
\]
where $\|(F,G)\| = \|F\| + \|G\|$ and $\|F\| = \sum |a_i|$ for
$F(z,t) = \sum a_i z^i t^{d-i}$.  Since coefficients are bounded and
evaluation on the unit disc satisfies \eqref{eq:eval_bound} below, the
Hopf endomorphism $\RS$ maps $S^3$ into $S^3$ without overflow.
Specifically, for $(z,t) \in D^2 \times D^2 := \{|z|\le 1\}\times\{|t|\le
1\}$,
\begin{equation}\label{eq:eval_bound}
  |F(z,t)| \le \|F\|, \qquad |G(z,t)| \le \|G\|,
\end{equation}
so all arithmetic stays within a bounded domain regardless of the orbit
length \cite[Prop.~2]{Alvarez2023}.

\subsection{The augmented projective line and the super-zero}

The map $f = A/B$ is undefined at points where both $A(z_0) = 0$ and
$B(z_0) = 0$ simultaneously, i.e.~at the \emph{super-zero} point
$[0:0] \notin \Pone$.  Following \cite[Def.~9]{Alvarez2023}, we work
in the \emph{augmented projective line}
$\Poneplus = \Pone \cup \{\tilde{0}\}$,
where $\tilde{0}$ is the image of the origin under the quotient map.  When
the r-irreducible pair $(F,G)$ satisfies $G(z_0,t_0) \ne 0$ for
$(z_0,t_0) \in S^3$ (Proposition~3(ii) of \cite{Alvarez2023}), the map
$\RS$ never reaches $\tilde{0}$ and remains well-defined on all of $S^3$.

\subsection{The spherical derivative in homogeneous coordinates}

For a rational map $f = A/B$ with r-irreducible pair $(F,G)$, the
\emph{spherical derivative} $\fsh: \Chat \to \R_+$ can be written as
\[
  \fsh(x) = \left|\frac{U(z,t)}{V(z,t)}\right|, \qquad x = z/t,
\]
where $U, V: \C^2 \to \C$ are the bihomogeneous polynomials of
\cite[Def.~16]{Alvarez2023}, continuous on all of $\C^2$ with $V(z,t)\ne
0$ for $(z,t)\ne(0,0)$ (Proposition~3(ii)).  This gives a continuous
map $\phi^S: S^3 \to \R_+$ satisfying $\phi^S(z,t) = \fsh(q(z,t))$,
well defined even at poles and at $\infty$.

\subsection{Lyapunov approximations}

Given a continuous map $\phi: X \to \R_+$ on a metric space $(X,g)$ and
$r \le s$ in $\N$, define
\[
  L_g^{[r,s]}(\phi)(x)
  = \left(\prod_{k=r}^{s} \phi(g^k(x))\right)^{1/(s+1-r)},
  \qquad x \in X.
\]
The \emph{Lyapunov function} $L_g(\phi)(x) = \lim_{s\to\infty}
L_g^{[0,s]}(\phi)(x)$ is constant on basins of attracting cycles:
$L_f(\phi)|_{B_f(C_i)} = c_i$ (almost everywhere), where the \emph{Lyapunov
constant} $c_i$ of cycle $C_i$ equals the $q_i$-th geometric mean of the
spherical derivative around the cycle \cite[Prop.~5]{Alvarez2023}.

%% ===========================================================================
\section{Sullivan's classification of Fatou components}
\label{sec:sullivan}
%% ===========================================================================

We recall the classification theorem that motivates our work.

\begin{theorem}[Sullivan \cite{Sullivan1983}]
  Let $f:\Chat\to\Chat$ be a rational map of degree $d\ge 2$.  Then $f$
  has no wandering Fatou components.  Every periodic Fatou component is of
  one of the following types:
  \begin{enumerate}[label=(\roman*)]
    \item \emph{Attracting basin}: $f^q$ has an attracting fixed point
          $\zeta$ inside with multiplier $|\lambda|<1$.
    \item \emph{Superattracting basin}: same with $\lambda=0$
          (i.e.~$\zeta$ is a critical point of $f^q$).
    \item \emph{Parabolic basin}: $f^q$ has a fixed point $\zeta$ on the
          boundary with multiplier $\lambda=e^{2\pi ip/q}$ a root of unity.
    \item \emph{Siegel disc}: $f^q$ is conformally conjugate to an
          irrational rotation $z\mapsto e^{2\pi i\theta}z$ on a disc; the
          fixed point $\zeta$ on the boundary has multiplier
          $\lambda=e^{2\pi i\theta}$ with $\theta\in\R\setminus\Q$
          satisfying the Brjuno condition.
    \item \emph{Herman ring}: $f^q$ is conformally conjugate to an
          irrational rotation on an annulus (no fixed point in the interior;
          occurs only for maps of degree $\ge 2$).
  \end{enumerate}
\end{theorem}

In cases (i) and (ii) the Lyapunov constant $c_i < 1$ (resp.~$c_i = 0$)
and the orbit converges geometrically (resp.~super-geometrically) to
$\zeta$.  Cases (iii)--(v) all have Lyapunov exponent $0$, i.e.~Lyapunov
constant $c_i = 1$, but differ in the \emph{rate and structure} of the
orbit behaviour near $\zeta$ (or the invariant annulus).

\begin{remark}\label{rem:previous}
  In \cite{Alvarez2023}, examples were chosen so that all Fatou components
  were attracting basins and the Julia set was the complement of their
  union.  The algorithms there detect only cases~(i) and (ii).  The present
  paper introduces criteria to handle cases~(iii)--(v).
\end{remark}

%% ===========================================================================
\section{Lyapunov-based criteria for attracting and parabolic basins}
\label{sec:lyapunov}
%% ===========================================================================

\subsection{Attracting and superattracting basins (recap)}

We briefly recall the criterion proved in \cite{Alvarez2023}.

\begin{proposition}[\cite{Alvarez2023}, Prop.~5--7]\label{prop:attracting}
  Let $C = (x_0, f(x_0), \ldots, f^{q-1}(x_0))$ be a $q$-cycle of $f$
  with Lyapunov constant
  $c = \bigl(\phi(x_0)\phi(f(x_0))\cdots\phi(f^{q-1}(x_0))\bigr)^{1/q}$
  where $\phi = \fsh$.
  \begin{enumerate}[label=(\roman*)]
    \item $C$ is attracting iff $c < 1$; superattracting iff $c = 0$.
    \item For every $y \in q^{-1}(B_f(C)) \subset S^3$, one has
          $L_{f^S}(\phi^S)(y) = c$.
    \item The approximation sequence $L_{f^S}^{[0,k]}(\phi^S)(y)$
          satisfies $|L_{f^S}^{[0,k-1]}(\phi^S)(y) -
          L_{f^S}^{[0,k]}(\phi^S)(y)| \leq 10^{-C}$
          for all $k \ge k_{(C,r)}$ (geometric convergence at rate $c$).
  \end{enumerate}
\end{proposition}

The convergence rate in part~(iii) is geometric: the gap between
consecutive approximations decays as $O(c^k)$ with $c < 1$.

\subsection{Parabolic basins: polynomial convergence}

Let $f^q$ have a parabolic fixed point $\zeta$ with multiplier
$\lambda = e^{2\pi ip/q}$ and parabolic multiplicity $m \ge 1$ (so
$f^q(z) = \zeta + (z-\zeta) + a_m(z-\zeta)^{m+1} + O((z-\zeta)^{m+2})$
with $a_m \ne 0$).  By the Leau--Fatou flower theorem \cite{Milnor2006},
there are $m$ attracting petals adjacent to $\zeta$, and every orbit in the
parabolic basin of $\zeta$ converges to $\zeta$ along one of these petals.

The key estimate is:

\begin{proposition}[Parabolic convergence rate]\label{prop:parabolic}
  Let $\zeta$ be a parabolic fixed point of $f$ with multiplicity $m$ and
  let $x \in B_f^{\mathrm{par}}(\zeta)$ (the parabolic basin).  Then:
  \begin{enumerate}[label=(\roman*)]
    \item $L_f(\phi)(x) = 1$ (the Lyapunov constant equals $1$).
    \item The orbit $\{f^n(x)\}_{n\ge 0}$ satisfies
          $|f^n(x) - \zeta| \sim (m n)^{-1/m}$ as $n\to\infty$
          (Parabolic Fatou coordinate estimate \cite{Milnor2006}).
    \item Consequently, the spherical derivative satisfies
          $\phi^S((\RS)^n(y)) = \fsh(f^n(x)) = 1 - c\,|f^n(x)-\zeta|^m +
          O(|f^n(x)-\zeta|^{m+1})$ for $y = q^{-1}(x) \in S^3$.
    \item The Lyapunov approximations satisfy
          \begin{equation}\label{eq:parabolic_rate}
            1 - L_{f^S}^{[0,N]}(\phi^S)(y)
            \sim \frac{C_m}{\log N}
            \quad\text{as }N\to\infty,
          \end{equation}
          for a constant $C_m > 0$ depending only on $m$ and the local
          dynamics near $\zeta$.
  \end{enumerate}
\end{proposition}

\begin{remark}
  Equation~\eqref{eq:parabolic_rate} implies that for a fixed large
  tolerance $10^{-C}$, the algorithm of \cite{Alvarez2023} will not
  declare convergence for orbits in the parabolic basin: the gap
  $|L_{f^S}^{[0,N-1]} - L_{f^S}^{[0,N]}|$ decreases to zero, but
  much more slowly than in the attracting case.  This is precisely the
  observable we exploit to detect parabolic basins.
\end{remark}

\begin{conjecture}\label{conj:parabolic_hopf}
  Proposition~\ref{prop:parabolic}(iii) holds when the spherical
  derivative $\phi^S$ is computed via the homogeneous bihomogeneous pair
  $(U,V)$ of \cite[Def.~16]{Alvarez2023}, i.e.~the normalisation to $S^3$
  does not distort the parabolic rate estimate.
\end{conjecture}

\noindent\textit{Evidence.}  Since $\RS: S^3 \to S^3$ is a smooth map
homotopic to $q^{-1} \circ f \circ q$ (Proposition~2 of \cite{Alvarez2023})
and the map $q: S^3 \to \Pone$ is a local isometry up to a constant
factor near any non-critical point, the local dynamics near a parabolic
point $\zeta$ is reproduced faithfully on $S^3$ except near the Hopf
fibres over $\zeta$.  A full proof requires a careful analysis of the
interaction between the normalisation and the Fatou coordinate, which we
intend to carry out in the final version.

\subsection{Practical criterion: detecting the convergence regime}

The following algorithm uses the asymptotic behaviour to distinguish the
three cases:

\begin{algorithm_desc}[Lyapunov regime detection]\label{alg:lyapunov}
  Input: a point $y \in S^3$, a maximum number of iterations $M$,
  a tolerance $\epsilon > 0$, and a polynomial-rate threshold $\alpha_0$.
  \begin{enumerate}[label=Step \arabic*.]
    \item Compute the sequence $a_k := L_{f^S}^{[0,k]}(\phi^S)(y)$ for
          $k = 0, 1, \ldots, M$.
    \item If $a_M < 1 - \epsilon$ and
          $|a_{M-1} - a_M| \le \epsilon^2$:
          declare \textsc{attracting} basin and record $c \approx a_M$.
    \item Else, if $|a_M - 1| \le \epsilon$ and
          $1 - a_k \ge C k^{-\alpha_0}$ for some $C > 0$ fitted by least
          squares on $\{(k, 1-a_k)\}$:
          declare \textsc{parabolic} basin.
    \item Else: defer to the rotation-number algorithm
          (Algorithm~\ref{alg:rotation}).
  \end{enumerate}
\end{algorithm_desc}

%% ===========================================================================
\section{Rotation-number algorithm for Siegel discs and Herman rings}
\label{sec:rotation}
%% ===========================================================================

\subsection{Rotation numbers on \texorpdfstring{$S^3$}{S3}}

Inside a Siegel disc or Herman ring, the map $f^q$ is analytically
conjugate to a rigid rotation $R_\theta: z \mapsto e^{2\pi i\theta} z$ on
a disc or annulus.  Every orbit $\{f^{qn}(x)\}_{n\ge 0}$ is quasi-periodic
with frequency $\theta$ (the rotation number).  The Lyapunov exponent is
exactly $0$, so $\fsh(f^n(x)) \to 1$ along the orbit but without
monotone convergence.

Working on $S^3$, let $y \in S^3$ with $x = q(y) \in \Pone$ lying in a
Siegel disc or Herman ring.  Define the \emph{angular sequence}
\begin{equation}\label{eq:angle_seq}
  \theta_n := \frac{1}{2\pi}\Arg\!\left(\frac{\langle (\RS)^n(y),\, y\rangle}
  {|\langle (\RS)^n(y),\, y\rangle|}\right), \qquad n = 1, 2, \ldots,
\end{equation}
where $\langle\cdot,\cdot\rangle$ denotes the standard Hermitian inner
product on $\C^2$.  When the orbit lies in a Siegel disc with rotation
$e^{2\pi i\theta}$, the sequence $\{n\theta \pmod 1\}$ approximates
$\{\theta_n\}$, and the rotation number can be recovered by any of the
following methods:
\begin{enumerate}[label=(\alph*)]
  \item \emph{Arithmetic mean}: $\theta \approx \frac{1}{N}\sum_{n=1}^N
        \theta_n / n$ (slow but robust).
  \item \emph{Discrete Fourier transform}: the dominant frequency in the
        DFT of the sequence $\{e^{2\pi i \theta_n}\}$ gives $\theta$.
  \item \emph{Continued fractions}: the Stern--Brocot algorithm applied to
        $\theta_N := \theta_N / N$ gives rational approximations
        $p_k/q_k \to \theta$, allowing early termination with a certified
        bound \cite{Yoccoz1988}.
\end{enumerate}

\begin{remark}[Herman rings]
  A Herman ring has no fixed point in its interior, so the inner product
  $\langle (\RS)^n(y), y\rangle$ may not accumulate near $y$.  One instead
  looks for the smallest $q$ such that the orbit $\{(\RS)^{qn}(y)\}$ appears
  to converge quasi-periodically, and applies \eqref{eq:angle_seq} to
  the $q$-th iterate.
\end{remark}

\subsection{Detecting rotation domains vs.~the Julia set}

A point in the Julia set also fails the attracting criterion and may also
fail the parabolic criterion.  The key distinction is that Julia set
orbits are \emph{not} quasi-periodic: they exhibit sensitive dependence
on initial conditions, while rotation-domain orbits are stable.

\begin{proposition}[Lyapunov exponent in rotation domains]
  \label{prop:siegel}
  Let $U$ be a Siegel disc or Herman ring of $f$.  For every $y \in
  q^{-1}(U) \subset S^3$:
  \begin{enumerate}[label=(\roman*)]
    \item $L_{f^S}(\phi^S)(y) = 1$.
    \item The sequence $\phi^S((\RS)^k(y))$ is quasi-periodic: for every
          $\varepsilon > 0$ there exists $N(\varepsilon)$ such that
          $|\phi^S((\RS)^{k+N(\varepsilon)}(y)) -
          \phi^S((\RS)^k(y))| < \varepsilon$ for all $k\ge 0$.
    \item The map $k \mapsto \phi^S((\RS)^k(y))$ does not converge: its
          lim-sup and lim-inf are both~$1$ but they are approached
          from both sides.
  \end{enumerate}
\end{proposition}

\begin{proof}[Sketch]
  Part~(i) follows because $f^{\#}(f^n(x)) = 1$ for all $n$ and $x$ in a
  rotation domain (rigid rotation preserves the spherical metric up to a
  constant factor at each point, and after one period $q$ the product
  $\prod_{k=0}^{q-1} f^{\#}(f^k(x)) = 1$).  Part~(ii) follows from the
  quasi-periodicity of the orbit itself together with the continuity of
  $\phi^S$.  Part~(iii) follows because the orbit visits all points of the
  Hopf fibre $q^{-1}(x)$ densely (by ergodicity of the irrational
  rotation), so the spherical derivative oscillates.
\end{proof}

\begin{algorithm_desc}[Rotation-number detection]\label{alg:rotation}
  Input: a point $y\in S^3$, a maximum period $Q$, iteration count $N$.
  \begin{enumerate}[label=Step \arabic*.]
    \item For $q = 1, 2, \ldots, Q$: compute
          $\delta_q := \|(\RS)^{qN/Q}(y) - y\|_{S^3}$.
          If $\delta_q < \varepsilon$ for some $q$: candidate period found.
    \item Compute the angular sequence $\{\theta_n\}_{n=1}^N$ of
          \eqref{eq:angle_seq} with $\RS$ replaced by $(\RS)^q$.
    \item Apply the DFT to $\{e^{2\pi i\theta_n}\}$; extract dominant
          frequency $\hat\theta$.
    \item Check stability: compare $\hat\theta$ for two initial
          conditions $y, y'$ close in $S^3$.  If $|\hat\theta(y) -
          \hat\theta(y')| < \varepsilon$: declare \textsc{siegel} disc or
          \textsc{herman} ring with rotation number $\hat\theta$.
    \item If no quasi-periodic pattern is found: declare \textsc{julia}.
  \end{enumerate}
\end{algorithm_desc}

%% ===========================================================================
\section{The unified classification algorithm}
\label{sec:algorithm}
%% ===========================================================================

Combining Algorithms~\ref{alg:lyapunov} and~\ref{alg:rotation}, we obtain
a complete classification procedure.  The input is a rational map given as
an r-irreducible pair $(F,G) \in \C_h[z,t]_{\mathrm{ri}}\C_h[z,t]$.

\begin{algorithm_desc}[Complete Fatou component classifier]
  \label{alg:full}
  \quad\\
  \textbf{Phase 1: grid and Hopf normalisation.}
  Choose a grid $D \subset \Pone$ of test points.  For each $x=[z:t]\in D$
  with $|z|+|t|=1$, let $y=(z,t)\in S^3$.

  \textbf{Phase 2: attracting cycle detection.}
  Apply Algorithm~\ref{alg:lyapunov} Steps~1--2.  Collect all attracting
  Lyapunov constants $c_i < 1$ found, together with representative cycles
  $C_i$ (as in \cite[\S4.1]{Alvarez2023}).

  \textbf{Phase 3: remaining points.}
  For each $y$ not classified as \textsc{attracting} in Phase~2:
  \begin{enumerate}[label=3.\arabic*.]
    \item Apply Step~3 of Algorithm~\ref{alg:lyapunov} to test for
          \textsc{parabolic}.
    \item If not parabolic, apply Algorithm~\ref{alg:rotation}.
    \item If Algorithm~\ref{alg:rotation} returns \textsc{julia}:
          the point is in $J(f)$ (or in the backward orbit of a critical
          point).
  \end{enumerate}

  \textbf{Phase 4: colouring.}
  Assign colours as follows:
  \begin{itemize}
    \item Each attracting basin $B_f(C_i)$: a distinct colour, graded by
          the band index $\lfloor \log_b L_{f^S}^{[r,s]}(\phi^S) \rfloor$
          (as in \cite{Alvarez2023}).
    \item Parabolic basins: a separate colour scheme, graded by the
          convergence progress $1 - L^{[0,N]}$.
    \item Siegel disc / Herman ring: a uniform hue, with saturation
          encoding the rotation number $\hat\theta$.
    \item Julia set: black (or by an escape-time / Böttcher coordinate
          where applicable).
  \end{itemize}
\end{algorithm_desc}

\begin{remark}[Complexity]
  The dominant cost is Phase~2 (attracting basin detection), already
  present in \cite{Alvarez2023}: $O(M)$ evaluations of $\RS$ and $\phi^S$
  per grid point.  Phase~3 adds at most $O(N)$ extra evaluations for the
  unclassified points; in practice these form a small fraction of the grid
  since they lie near $J(f)$.  The total cost is comparable to that of the
  original algorithm.
\end{remark}

%% ===========================================================================
\section{Examples and applications}
\label{sec:examples}
%% ===========================================================================

\subsection{A rational map with a parabolic fixed point: $f(z) = z^2 + \tfrac{1}{4}$}

The map $f(z) = z^2 + \frac{1}{4}$ has a parabolic fixed point at
$\zeta = \frac{1}{2}$ with multiplier $\lambda = f'(\frac{1}{2}) = 1$
and multiplicity $m=1$ (one attracting petal, pointing toward $\zeta$
from the left).  The corresponding homogeneous pair is
\[
  F(z,t) = z^2 + \tfrac{1}{4}t^2, \qquad G(z,t) = t^2.
\]
The normalised Hopf endomorphism is
\[
  \RS(z,t) = \frac{(z^2 + \tfrac{1}{4}t^2,\, t^2)}{|z^2 +
  \tfrac{1}{4}t^2| + |t^2|}.
\]
The Julia set $J(f)$ is the boundary of the filled-in Julia set, which is
connected (since the unique critical point $z=0$ is on the boundary of
the parabolic basin).  Every point with bounded orbit converges to
$\zeta = \frac{1}{2}$ along the unique attracting petal.

In this example Algorithm~\ref{alg:lyapunov} should classify the entire
basin as \textsc{parabolic} with $m = 1$, and $J(f)$ as \textsc{julia}.
Since there are no attracting cycles, Phase~2 returns no Lyapunov
constants below $1$, and all bounded orbits proceed to Phase~3.1.

\subsection{A quadratic map with a Siegel disc}
\label{sec:example_siegel}

Let $\theta = (\sqrt{5}-1)/2$ (the golden mean, which satisfies the
Brjuno condition).  By the theory of quadratic polynomials
\cite{Siegel1942, Yoccoz1988}, the map
\[
  f_\theta(z) = e^{2\pi i\theta} z + z^2
\]
has a Siegel disc $U_\theta$ centred at $0$ with rotation number
$\theta$.  The boundary $\partial U_\theta$ is a quasi-circle and
$J(f_\theta) = \partial U_\theta$.

The corresponding homogeneous pair is
\[
  F(z,t) = e^{2\pi i\theta} zt + z^2, \qquad G(z,t) = t^2.
\]
For $y \in q^{-1}(U_\theta)$, the angular sequence \eqref{eq:angle_seq}
should give $\hat\theta = (\sqrt{5}-1)/2$, detectable via the DFT by
the dominant pair of peaks at $\pm\hat\theta$.

\begin{remark}
  The parameter $\alpha = e^{2\pi i\theta}$ is a point on the boundary of
  the main cardioid of the Mandelbrot set (where the map has an
  indifferent fixed point), but strictly speaking the Mandelbrot set
  question (whether the Siegel disc touches $J(f)$) is separate.  For
  our classification purposes, any Brjuno $\theta$ gives an example.
\end{remark}

\subsection{Newton's method: detecting parameter-space $n$-cycles}

This extends \cite[\S6]{Alvarez2023}, which applied the Lyapunov
framework to Newton's method for cubic polynomials of the form
$p_\alpha(x) = (x-\alpha)(x^2+1)$.  For values of $\alpha$ in the
Mandelbrot-like set of Figure~1 of \cite{Alvarez2023}, Newton's method
has attracting $n$-cycles for $n \ge 2$ rather than converging to a root.

With the extended algorithm, one can further distinguish:
\begin{itemize}
  \item \emph{Red regions} in \cite[Fig.~2--3]{Alvarez2023}: orbits where
        no attracting cycle was found.  Among these, our algorithm
        identifies the subset corresponding to parabolic parameters
        (boundary of the Mandelbrot-like bulbs) and to Siegel parameters
        (points of the Mandelbrot set with Brjuno rotation).
  \item The \emph{bifurcation locus} is precisely the boundary between
        attracting regions of different periods, and should correspond to
        parabolic parameters where the Lyapunov constant equals $1$.
\end{itemize}
This provides a finer picture of the parameter space than the binary
``converges / does not converge'' output of the original algorithm.

\subsection{Halley's method and other rational iteration schemes}

The homogeneous-pair formalism applies equally to Halley's method and
any iterative method representable as a rational map.  For Halley's
method applied to $p(x) = x^n - 1$, the basins and their boundaries
are qualitatively similar to Newton's but with higher symmetry.  The
same classification algorithm applies without modification, making
\texttt{RiemannBasins.jl} a unified platform for all such methods.

%% ===========================================================================
\section{Implementation notes}
\label{sec:implementation}
%% ===========================================================================

The algorithms described above extend the \texttt{LCD.jl}
(Lyapunov Cycle Detector) package \cite{LCD2022} and the
\texttt{RiemannBasins.jl} package \cite{RiemannBasins}.  The
implementation uses the Julia programming language \cite{Julia2017},
which is well suited to this task due to its high numerical performance,
its native support for complex arithmetic, and the existing
\texttt{FFTW.jl} library for the DFT-based rotation number estimation.

The key new exported functions planned for the extended package are:
\begin{itemize}
  \item \texttt{classify\_fatou(f, grid; tol, max\_iter)} --- returns a
        label array with entries from
        \texttt{\{:attracting, :superattracting, :parabolic, :siegel,
        :herman, :julia\}}.
  \item \texttt{rotation\_number(f, z; max\_iter)} --- estimates the
        rotation number at a point in a potential rotation domain.
  \item \texttt{parabolic\_portrait(f; ...)} --- full portrait with the
        parabolic basin coloured separately.
\end{itemize}

Basic functions used to write the new code had been previously developed
in \cite{Extremiana2019} and \cite{PBURF2019}, following the tradition
established in \cite{HernandezParicio2012}.

%% ===========================================================================
\section{Conclusions and future work}
\label{sec:conclusions}
%% ===========================================================================

We have presented the theoretical foundations and algorithmic design for
a complete classification of Fatou components of rational maps within the
Hopf fibration framework of \cite{Alvarez2023}.  Working entirely in
normalised homogeneous coordinates on $S^3$ avoids the numerical
pathologies (overflow, underflow, indeterminacy at $[0:0]$) that affect
the standard affine approach.

The two main new technical contributions are:
\begin{enumerate}[label=(\arabic*)]
  \item The \emph{parabolic detection criterion} based on the polynomial
        convergence rate of the Lyapunov approximation sequence, which
        allows one to distinguish parabolic basins from attracting ones
        and from the Julia set without computing Fatou coordinates
        explicitly.
  \item The \emph{rotation-number algorithm on $S^3$}, which detects
        Siegel discs and Herman rings and returns the rotation number as
        a certified rational approximation.
\end{enumerate}

Several directions for future work remain open:

\begin{itemize}
  \item \textbf{Proof of Conjecture~\ref{conj:parabolic_hopf}.}  A
        full proof that the normalisation to $S^3$ does not distort the
        parabolic rate estimate requires a detailed analysis of the
        interaction between the Hopf fibration and the Fatou coordinate
        near a parabolic point.

  \item \textbf{Hausdorff dimension estimates.}  The techniques of
        \cite{Alvarez2023} can potentially be adapted to estimate the
        Hausdorff dimension of $J(f)$ via the exponent $\delta(f)$ of
        Sullivan's $\delta$-conformal measure \cite{Sullivan1983,
        Gelfert2010}.  The classification of Fatou components is a
        prerequisite for such estimates in the non-hyperbolic case.

  \item \textbf{Parameter-space applications.}  As illustrated in
        Section~\ref{sec:examples}, the algorithm provides a finer
        decomposition of parameter spaces (the Mandelbrot-like sets of
        \cite{Alvarez2023}) by distinguishing parabolic parameters on
        bifurcation boundaries from Siegel parameters inside bulbs.
        This connects to the program of \cite{McMullen1987} on the
        universality of the Mandelbrot set.

  \item \textbf{Semigroup actions.}  The framework extends naturally to
        the iteration of rational semigroups generated by two or more
        rational maps, as studied in \cite{HinkkaneMartin} (reference
        to be added).  The Lyapunov function $L_g(\phi^S)$ depends on
        the choice of generator $g$, and understanding how the Fatou
        set of the semigroup decomposes with respect to individual
        generators is an open problem.

  \item \textbf{Visualisation of bifurcation loci.}  The parabolic
        detection algorithm provides a numerical means to trace
        bifurcation loci in parameter space, which are the boundaries
        between regions with different attracting periods.  These
        boundaries have fractal structure and their numerical
        visualisation is currently limited by the speed of convergence
        near parabolic parameters.
\end{itemize}

\bigskip
\noindent\textbf{Acknowledgements.}
The authors thank Professor J.M.~Rodríguez Sanjurjo for his continued
interest and encouragement.  This work is a continuation of the program
initiated in \cite{Extremiana2019,Alvarez2023}.  Computations were
performed using the Julia language \cite{Julia2017} with the
\texttt{RiemannBasins.jl} package \cite{RiemannBasins}.
L.J.H.-P.\ acknowledges support from project REGI25/61 of the
Universidad de La Rioja.

%% ===========================================================================
\bibliographystyle{plain}
\bibliography{draft_fatou_hopf}
%% ===========================================================================

\end{document}
